\section{A typed shift-reduce implementation}
The parser uses the invariant that the state stack contains exactly one more element than the semantic-value stack.
A table entry contains \texttt{Shift state}, \texttt{Reduce rule}, or \texttt{Accept}.
The top of an F\# list is its head, so pushing a value uses \texttt{::} rather than appending at the end.
Reduction removes exactly the production body's length and applies the goto transition to the exposed state.
An empty production removes zero entries, which is handled explicitly by \texttt{List.take 0} and \texttt{List.skip 0}.
@@snippet:shift-reduce@@

\texttt{parseBy} accepts a leaf function and a reduction function. The recognizer supplies unit values;
the AST parser supplies tree constructors. This is a small real interface between table-driven parsing and semantic actions.
It keeps the actual transition algorithm in one place. An absent entry, an unresolved conflict or an invalid stack
produces \texttt{Error message}; acceptance produces \texttt{Ok value}.

The abstract action and goto tables below remain the mathematical example from the original lecture.
Their state numbering is illustrative; generated F\# table state numbers follow structural ordering.
